\documentclass[12pt]{article}
\usepackage[utf8]{inputenc}
\usepackage[T1]{fontenc}
\usepackage[margin=1in]{geometry}
\usepackage{setspace}
\usepackage{xcolor}
\usepackage[colorlinks=true, urlcolor=blue!60!black, linkcolor=blue!60!black]{hyperref}
\setlength{\parindent}{2em}
\setlength{\parskip}{1.25em}
\renewcommand{\baselinestretch}{1.0}
\begin{document}
\begin{center}
{\LARGE \textit{Timbs v.\ Indiana}, 586 U.S. 146 (2019)\par}
\vspace{0.8em}
{\large Case Brief\par}
\end{center}
\vspace{1em}

\section{Memory Jogger}

The Supreme Court unanimously held that the Eighth Amendment's Excessive Fines Clause is incorporated against the states through the Fourteenth Amendment, barring Indiana from claiming immunity from federal excessive-fines scrutiny when seeking civil \textit{in rem} forfeiture of a \$42,000 Land Rover for a drug offense carrying a \$10,000 maximum fine (\href{https://www.courtlistener.com/opinion/4591916/timbs-v-indiana/}{CourtListener}; \href{https://www.law.cornell.edu/supremecourt/text/17-1091}{LII}).

\section{Detailed Case Facts}

After his father's death, Tyson Timbs received life insurance proceeds and used approximately \$42,000 to purchase a 2012 Land Rover LR2 SUV (\href{https://www.law.cornell.edu/supremecourt/text/17-1091}{586 U.S. at 148}; \href{https://caselaw.findlaw.com/court/in-supreme-court/2028477.html}{FindLaw}). In May 2013, while struggling with an opioid addiction, Timbs drove the Land Rover in Grant County, Indiana, to sell small amounts of heroin (about four grams total) to undercover officers (\href{https://ij.org/press-release/14-months-after-u-s-supreme-court-victory-tyson-timbs-wins-civil-forfeiture-case/}{Institute for Justice}; \href{https://caselaw.findlaw.com/court/in-supreme-court/2028477.html}{FindLaw}). Police arrested Timbs on May 31, 2013, and seized the Land Rover (\href{https://caselaw.findlaw.com/court/in-supreme-court/2028477.html}{FindLaw}).

Timbs pleaded guilty in Indiana state court to one count of dealing in a controlled substance and one count of conspiracy to commit theft (\href{https://www.law.cornell.edu/supremecourt/text/17-1091}{LII}). The trial court sentenced him to six years (five suspended), consisting of one year of home detention and five years of probation including a court-supervised addiction-treatment program, and ordered him to pay \$1,203 in fees and costs (\href{https://www.law.cornell.edu/supremecourt/text/17-1091}{LII}; \href{https://fedsoc.org/case/timbs-v-indiana}{Federalist Society}). Meanwhile, the State of Indiana hired a private law firm to bring a civil \textit{in rem} forfeiture action against the Land Rover under Indiana Code \S~34-24-1-1, alleging the vehicle had been used to transport heroin (\href{https://www.law.cornell.edu/supremecourt/text/17-1091}{LII}; \href{https://caselaw.findlaw.com/court/in-supreme-court/2028477.html}{FindLaw}). Because the maximum statutory monetary fine for Timbs's felony drug conviction was \$10,000, the value of his legally purchased Land Rover was more than four times the statutory maximum fine (\href{https://www.law.cornell.edu/supremecourt/text/17-1091}{LII}).

\section{Procedural History}

Following a hearing on the civil forfeiture complaint, the Grant County Superior Court found that the Land Rover had been used to facilitate a criminal violation, but denied forfeiture because seizing a \$42,000 vehicle---more than four times the \$10,000 statutory maximum fine---was ``grossly disproportionate to the gravity of Timbs's offense'' and therefore unconstitutional under the Eighth Amendment's Excessive Fines Clause (\href{https://www.law.cornell.edu/supremecourt/text/17-1091}{LII}). The Court of Appeals of Indiana affirmed, 62 N.E.3d 472 (Ind. Ct. App. 2016) (\href{https://en.wikipedia.org/wiki/Timbs_v._Indiana}{Wikipedia}).

The Indiana Supreme Court reversed, 84 N.E.3d 1179 (Ind. 2017), declining to decide whether the forfeiture was excessive and instead holding that the U.S. Supreme Court had never incorporated the Excessive Fines Clause against the states, so the Clause constrained only federal action (\href{https://www.law.cornell.edu/supremecourt/text/17-1091}{LII}). Represented by the Institute for Justice, Timbs petitioned for certiorari, which the U.S. Supreme Court granted on June 18, 2018 (585 U.S. 1002) (\href{https://www.govinfo.gov/content/pkg/USREPORTS-586/pdf/USREPORTS-586-146.pdf}{U.S. Reports}; \href{https://www.scotusblog.com/cases/timbs-v-indiana/}{SCOTUSblog}). The case was argued on November 28, 2018, and decided on February 20, 2019, vacating and remanding to the Indiana Supreme Court (which ultimately held the forfeiture unconstitutionally excessive on remand) (\href{https://www.law.cornell.edu/supremecourt/text/17-1091}{LII}; \href{https://en.wikipedia.org/wiki/Timbs_v._Indiana}{Wikipedia}).

\section{Judicial Votes}

\begin{itemize}
  \item \textbf{Vote:} 9--0, vacating the judgment of the Supreme Court of Indiana and remanding (\href{https://www.law.cornell.edu/supremecourt/text/17-1091}{LII}; \href{https://www.courtlistener.com/opinion/4591916/timbs-v-indiana/}{CourtListener}).
  \item \textbf{Majority (8):} Justice Ruth Bader Ginsburg (author), joined by Chief Justice John G. Roberts, Jr., and Justices Stephen G. Breyer, Samuel A. Alito, Jr., Sonia Sotomayor, Elena Kagan, Neil M. Gorsuch, and Brett M. Kavanaugh (\href{https://www.law.cornell.edu/supremecourt/text/17-1091}{LII}).
  \item \textbf{Concurrence (1):} Justice Neil M. Gorsuch joined the majority opinion in full and filed a brief separate concurring opinion (\href{https://www.law.cornell.edu/supremecourt/text/17-1091}{LII}).
  \item \textbf{Concurrence in the Judgment (1):} Justice Clarence Thomas filed an opinion concurring in the judgment only (\href{https://www.law.cornell.edu/supremecourt/text/17-1091}{LII}).
\end{itemize}

\section{Holding}

\textbf{Judgment:} Vacated and remanded. The Eighth Amendment's Excessive Fines Clause applies to the states (\href{https://www.law.cornell.edu/supremecourt/text/17-1091}{LII}).

\textbf{Rule Announced:} The Eighth Amendment's Excessive Fines Clause is an incorporated protection applicable to the states under the Fourteenth Amendment's Due Process Clause because protection against excessive punitive economic sanctions is both ``fundamental to our scheme of ordered liberty'' and ``deeply rooted in this Nation's history and tradition'' (\href{https://www.law.cornell.edu/supremecourt/text/17-1091}{LII}). Once a Bill of Rights guarantee is incorporated, it applies identically to both the federal government and the states---including its established application to civil \textit{in rem} forfeitures that are at least partially punitive under \textit{Austin v. United States}, 509 U.S. 602 (1993) (\href{https://www.law.cornell.edu/supremecourt/text/17-1091}{LII}).

\textbf{Holding vs.\ Dicta:} The \textit{holding} is twofold: (1) the Excessive Fines Clause is incorporated against the states via the Fourteenth Amendment's Due Process Clause; and (2) the incorporation inquiry asks whether the guaranteed \textit{right itself}---not each particular application, such as civil \textit{in rem} forfeiture---is fundamental or deeply rooted. The Court explicitly \textit{did not hold} whether forfeiting Timbs's Land Rover was actually excessive, nor did it establish a specific multi-factor excessiveness test or decide whether an offender's wealth or ability to pay is relevant (expressly noting that \textit{United States v. Bajakajian}, 524 U.S. 321, 340 n.15 (1998), left that question open). Its observations about modern state and local reliance on fines and fees for revenue, as well as its refusal to reconsider \textit{Austin} on procedural forfeiture grounds (neither pressed nor passed upon below), frame the context without deciding the merits of \textit{Austin} or of Timbs's specific forfeiture (\href{https://www.law.cornell.edu/supremecourt/text/17-1091}{LII}).

\section{Analysis of Opinions}

\subsection{Majority Opinion (Justice Ginsburg)}

Justice Ginsburg applied the modern selective-incorporation framework synthesized in \href{https://supreme.justia.com/cases/federal/us/561/742/}{\textit{McDonald v. City of Chicago}, 561 U.S. 742 (2010)}: a Bill of Rights guarantee is incorporated through the Fourteenth Amendment's Due Process Clause if it is ``fundamental to our scheme of ordered liberty'' or ``deeply rooted in this Nation's history and tradition,'' and once incorporated, ``there is no daylight between the federal and state conduct it prohibits or requires'' (\href{https://www.law.cornell.edu/supremecourt/text/17-1091}{LII}).

Tracing the Clause's ``venerable lineage,'' Ginsburg documented continuous Anglo-American protection against excessive fines from Magna Carta (1215)---which required amercements to be proportioned to the offense and to spare an offender's livelihood---through the English Bill of Rights (1689) enacted after Stuart abuses, the Virginia Declaration of Rights (1776), and state constitutions at the Founding (8 of 13 states, representing 70\% of the population in 1787), at the ratification of the Fourteenth Amendment (35 of 37 states, over 90\% of the population in 1868), and today (all 50 states) (\href{https://www.law.cornell.edu/supremecourt/text/17-1091}{LII}). She also underscored the Reconstruction history of Southern Black Codes, which used draconian vagrancy fines to coerce formerly enslaved people into involuntary labor, as well as the structural danger that---unlike imprisonment, which costs the state money---fines generate government revenue and can be wielded to chill political dissent (\href{https://www.law.cornell.edu/supremecourt/text/17-1091}{LII}).

Finally, Ginsburg rejected Indiana's two fallback arguments aimed at shielding civil \textit{in rem} forfeitures. First, Indiana's invitation to overrule \textit{Austin v. United States} (1993) was not properly before the Court because Indiana neither pressed that issue below nor had the right to expand the question presented. Second, Indiana misapprehended incorporation doctrine by arguing that \textit{in rem} forfeiture itself must be deeply rooted: just as the First Amendment applies to social media (\textit{Packingham v. North Carolina}) and the Fourth Amendment applies to cell phones (\textit{Riley v. California}) without re-litigating incorporation for each application, incorporation turns on the right itself, not every specific application (\href{https://www.law.cornell.edu/supremecourt/text/17-1091}{LII}).

\subsection{Concurring Opinion (Justice Gorsuch)}

Joining the majority opinion in full, Justice Gorsuch wrote a single-paragraph concurrence acknowledging that, ``[a]s an original matter,'' the ``appropriate vehicle for incorporation may well be the Fourteenth Amendment's Privileges or Immunities Clause, rather than, as this Court has long assumed, the Due Process Clause'' (\href{https://www.law.cornell.edu/supremecourt/text/17-1091}{LII}). Citing Justice Thomas's concurrences in \textit{Timbs} and \textit{McDonald} alongside scholarship by Bryan Wildenthal, Akhil Amar, and Michael Kent Curtis, Gorsuch nonetheless joined the Court's opinion because ``nothing in this case turns on that question'' (\href{https://www.law.cornell.edu/supremecourt/text/17-1091}{LII}).

\subsection{Concurring Opinion in the Judgment (Justice Thomas)}

Justice Thomas agreed that the Fourteenth Amendment binds states to the Excessive Fines Clause, but refused to join the majority's reliance on the Due Process Clause, condemning ``the oxymoronic `substantive' `due process' doctrine'' as a ``legal fiction'' lacking textual constraints---pointing to \textit{Dred Scott v. Sandford} and \textit{Roe v. Wade} as cautionary examples (\href{https://www.law.cornell.edu/supremecourt/text/17-1091}{LII}). Because Timbs did not claim Indiana denied him fair judicial procedures, his claim had nothing to do with ``process.''

Instead, Thomas grounded incorporation in the Privileges or Immunities Clause, maintaining that the ratifying public in 1868 understood ``privileges or immunities'' as synonyms for the inalienable, constitutionally enumerated rights of American citizenship (\href{https://www.law.cornell.edu/supremecourt/text/17-1091}{LII}). Through an extensive historical survey---from the Charter of Liberties of Henry I (1101), Magna Carta, the Star Chamber's abolition, and Stuart-era King's Bench abuses, through Anti-Federalist essays, Patrick Henry, Story, Kent, and 39th Congress debates over Black Code vagrancy fines---Thomas concluded that immunity from excessive fines was universally recognized as a privilege of American citizenship binding in full on the states (\href{https://www.law.cornell.edu/supremecourt/text/17-1091}{LII}).

\section{Examples: Future Applications}

\begin{enumerate}
  \item \textbf{Same-Side (Disproportionate State Civil Asset Forfeiture):} A state seizes and seeks civil \textit{in rem} forfeiture of a \$150,000 small airplane or family home because the owner committed a misdemeanor offense carrying a maximum \$2,500 fine. Under \textit{Timbs} and \textit{Austin}, the state cannot evade Eighth Amendment scrutiny by labeling the proceeding ``civil''; because the forfeiture is at least partially punitive, state courts must evaluate whether it is grossly disproportional to the gravity of the offense.
  \item \textbf{Same-Side (Municipal Code-Enforcement Fines):} A city imposes daily accruing fines totaling \$160,000 against a homeowner for minor property-maintenance infractions such as overgrown grass or a cracked driveway (\href{https://virginialawreview.org/wp-content/uploads/2025/03/Asbridge_Revised.pdf}{Virginia Law Review}). Under \textit{Timbs}, municipal punitive fines are directly constrained by the incorporated Eighth Amendment Excessive Fines Clause, preventing local governments from using exorbitant punitive tolls as a revenue-raising substitute for taxation.
  \item \textbf{Opposite-Side (Purely Remedial Civil Damages or Tax Collection):} A state sues a contractor to recover actual compensatory damages for defective highway work, or collects unpaid back taxes plus standard statutory interest strictly calibrated to reimburse the public treasury. Because the Excessive Fines Clause reaches only payments extracted ``as punishment for some offense'' (\href{https://www.law.cornell.edu/supremecourt/text/17-1091}{LII}), purely remedial assessments with no punitive purpose fall outside \textit{Timbs}.
  \item \textbf{Opposite-Side (Forfeiture of Direct Contraband or Criminal Proceeds):} State police confiscate \$50,000 in cash proven to be the direct proceeds of illegal narcotics trafficking, or seize illegal automatic weapons classified as contraband \textit{per se}. Even after \textit{Timbs}, forfeiting contraband or the direct illegal proceeds of crime is generally treated as remedial disgorgement rather than an excessive punitive fine.
  \item \textbf{Fence-Sitter (Ability-to-Pay Challenges to Modest Statutory Fines Against Indigent Defendants):} A state imposes a \$1,500 mandatory statutory fine---well within the statutory maximum---on an indigent defendant convicted of a low-level driving offense, triggering license suspension and debt collection that deprive the defendant of the ability to earn a livelihood. Although \textit{Timbs} repeatedly invoked Magna Carta's ``saving his contenement'' (livelihood) principle, the Court explicitly noted that \textit{Bajakajian} took no position on whether a defendant's income and wealth govern excessiveness (\href{https://www.law.cornell.edu/supremecourt/text/17-1091}{LII}; \href{https://harvardlawreview.org/print/vol-133/timbs-v-indiana/}{Harvard Law Review}), leaving lower courts divided on how much weight to give inability to pay.
\end{enumerate}

\section{Critique}

\textbf{Originalist and Textualist Perspectives:} Originalist and libertarian commentators warmly welcomed the unanimous incorporation of the Excessive Fines Clause as a long-overdue check on state civil asset forfeiture, while critiquing the majority's continued reliance on substantive due process (\href{https://www.cato.org/supreme-court-review/2018-2019/timbs-v-indiana-mere-constitutional-housekeeping-or-timely-revival}{Cato Supreme Court Review}; \href{https://fedsoc.org/case/timbs-v-indiana}{Federalist Society}). Scholars and jurists aligned with Justices Thomas and Gorsuch argue that routing incorporation through the Due Process Clause perpetuates the textual distortion created by \textit{The Slaughter-House Cases} and invites unmoored judicial policymaking, whereas restoring the Privileges or Immunities Clause would ground incorporation in the Fourteenth Amendment's original public meaning (\href{https://texaslawreview.org/a-distinction-with-a-difference/}{Texas Law Review}; \href{https://www.cato.org/supreme-court-review/2018-2019/timbs-v-indiana-mere-constitutional-housekeeping-or-timely-revival}{Cato Supreme Court Review}). At the same time, some observers note that \textit{Timbs} left the core federalism question of \textit{Austin}'s validity unresolved, potentially inviting future originalist challenges to whether historical \textit{in rem} forfeitures were understood as ``fines'' in 1791 or 1868 (\href{https://harvardlawreview.org/print/vol-133/timbs-v-indiana/}{Harvard Law Review}).

\textbf{Progressive and Doctrinal Perspectives:} Progressive and criminal-justice scholars praised \textit{Timbs} for recognizing both the Reconstruction history of Black Code fines and the modern reality that state and local governments increasingly rely on fines, fees, and forfeitures as regressive revenue streams (\href{https://harvardlawreview.org/print/vol-133/timbs-v-indiana/}{Harvard Law Review}). Yet critics such as Beth Colgan and the \textit{Harvard Law Review} warn that \textit{Timbs} may prove largely symbolic unless courts move beyond rigid comparisons between asset values and statutory maximums to incorporate defendants' ability to pay and livelihood impacts---because even modest fines can devastate low-income people while failing to deter the wealthy (\href{https://harvardlawreview.org/print/vol-133/timbs-v-indiana/}{Harvard Law Review}). Doctrinal commentators also point out that leaving lower courts without a uniform excessiveness standard has produced fragmented state-court tests (\href{https://virginialawreview.org/wp-content/uploads/2025/03/Asbridge_Revised.pdf}{Virginia Law Review}).

\textbf{Underlying Competing Values:} \textit{Timbs} balances \textit{state sovereignty and law-enforcement discretion} (allowing states to structure criminal penalties, civil forfeitures, and local revenue systems) against \textit{uniform national protection of individual property and liberty} (preventing governments from weaponizing revenue-generating financial sanctions to chill dissent or impose disproportionate punishment), alongside a methodological clash between \textit{stare decisis} in Due Process incorporation and \textit{originalist textual fidelity} to the Privileges or Immunities Clause.

\section{Key Quotations}

\begin{quote}
``Incorporated Bill of Rights guarantees are `enforced against the States under the Fourteenth Amendment according to the same standards that protect those personal rights against federal encroachment.' Thus, if a Bill of Rights protection is incorporated, there is no daylight between the federal and state conduct it prohibits or requires.''
\end{quote}
\noindent --- Justice Ruth Bader Ginsburg, for the Court, 586 U.S. at 150 (quoting \textit{McDonald v. Chicago}, 561 U.S. at 765) (\href{https://www.law.cornell.edu/supremecourt/text/17-1091}{LII}; \href{https://www.courtlistener.com/opinion/4591916/timbs-v-indiana/}{CourtListener}).

\begin{quote}
``For good reason, the protection against excessive fines has been a constant shield throughout Anglo-American history: Exorbitant tolls undermine other constitutional liberties.''
\end{quote}
\noindent --- Justice Ruth Bader Ginsburg, for the Court, 586 U.S. at 153--154 (\href{https://www.law.cornell.edu/supremecourt/text/17-1091}{LII}).

\begin{quote}
``In considering whether the Fourteenth Amendment incorporates a protection contained in the Bill of Rights, we ask whether the right guaranteed---not each and every particular application of that right---is fundamental or deeply rooted.''
\end{quote}
\noindent --- Justice Ruth Bader Ginsburg, for the Court, 586 U.S. at 156 (\href{https://www.law.cornell.edu/supremecourt/text/17-1091}{LII}).

\begin{quote}
``As an original matter, I acknowledge, the appropriate vehicle for incorporation may well be the Fourteenth Amendment's Privileges or Immunities Clause, rather than, as this Court has long assumed, the Due Process Clause. \ldots\ But nothing in this case turns on that question, and, regardless of the precise vehicle, there can be no serious doubt that the Fourteenth Amendment requires the States to respect the freedom from excessive fines enshrined in the Eighth Amendment.''
\end{quote}
\noindent --- Justice Neil M. Gorsuch, concurring, 586 U.S. at 157 (\href{https://www.law.cornell.edu/supremecourt/text/17-1091}{LII}).

\begin{quote}
``Instead of reading the Fourteenth Amendment's Due Process Clause to encompass a substantive right that has nothing to do with `process,' I would hold that the right to be free from excessive fines is one of the `privileges or immunities of citizens of the United States' protected by the Fourteenth Amendment.''
\end{quote}
\noindent --- Justice Clarence Thomas, concurring in the judgment, 586 U.S. at 157 (\href{https://www.law.cornell.edu/supremecourt/text/17-1091}{LII}).


\end{document}
